\chapter{Systementwurf}
\label{chap:systementwurf}

Die Anwendung basiert auf der Leitidee, dass fachliche Korrektheit und nachvollziehbare Interaktion gemeinsam gedacht werden müssen. Ein Scheduling-Visualizer soll nicht nur Ergebnisse berechnen, sondern auch verständlich machen, warum ein Prozess zu einem bestimmten Zeitpunkt läuft und ein anderer wartet. Der Systementwurf verbindet daher Architekturentscheidungen mit didaktischen Anforderungen.

\section{Entwurfsziele und Leitprinzipien}

Aus der Anforderungsanalyse ergeben sich vier zentrale Entwurfsziele:

\begin{itemize}
    \item \textbf{Fachliche Konsistenz:} Simulation, Ereignisverlauf und Kennzahlen müssen aus einer gemeinsamen Zustandsbasis abgeleitet werden.
    \item \textbf{Didaktische Nachvollziehbarkeit:} Zustandswechsel und Präemptionen müssen zeitlich verständlich dargestellt werden.
    \item \textbf{Deterministik:} Identische Eingaben sollen zu identischen Abläufen führen.
    \item \textbf{Erweiterbarkeit:} Neue Algorithmen und Darstellungsformen sollen ohne grundlegenden Umbau integrierbar sein.
\end{itemize}

Diese Ziele werden durch eine komponentenbasierte Struktur mit klaren Verantwortlichkeiten umgesetzt. Die Kopplung zwischen Simulationslogik und Darstellung bleibt bewusst gering, damit fachliche Änderungen kontrolliert möglich sind und die UI als Vermittlungsschicht wirken kann.

\section{Gesamtarchitektur und Schichtenmodell}

Das System ist in drei logisch getrennte Ebenen gegliedert:

\begin{itemize}
    \item \textbf{Szenario- und Konfigurationslogik:} Verwaltung von Taskdaten, Parametern und Laufkonfigurationen.
    \item \textbf{Simulationskern:} Deterministische Berechnung des zeitlichen Ablaufs je Scheduling-Verfahren.
    \item \textbf{Darstellung und Steuerung:} Visualisierung der Ergebnisse sowie Interaktion (Start, Pause, Schrittmodus, Zeitnavigation, Vergleich).
\end{itemize}

Die Steuerung ist dabei Teil der Darstellungsschicht, weil sie den didaktischen Ablauf unmittelbar beeinflusst: Nutzende greifen über die Oberfläche in die Wiedergabe ein und erhalten direkt visuelles Feedback auf denselben fachlichen Daten.

\begin{figure}[htbp]
    \centering
    \begin{tikzpicture}[
  font=\footnotesize,
  >=Latex,
  comp/.style={
    draw=black!65,
    rounded corners=2mm,
    very thick,
    minimum width=3.9cm,
    minimum height=1.15cm,
    align=center
  },
  layer/.style={
    draw=black!55,
    rounded corners=2mm,
    very thick,
    minimum width=13.4cm,
    minimum height=3.25cm
  },
  flow/.style={->, thick, draw=black!75},
  feedback/.style={->, thick, dashed, draw=black!65}
]

% Layer
\node[layer, fill=blue!5]   (L1) at (0,0) {};
\node[layer, fill=green!5]  (L2) at (0,-4.0) {};
\node[layer, fill=orange!6] (L3) at (0,-8.0) {};

% Titles
\node[anchor=west, font=\bfseries] at (-6.45,1.4) {Benutzeroberfläche (Darstellung \& Steuerung)};
\node[anchor=west, font=\bfseries] at (-6.45,-2.6) {Anwendungslogik};
\node[anchor=west, font=\bfseries] at (-6.45,-6.6) {Visualisierung \& Auswertung};

% Row 1
\node[comp, fill=blue!14] (ui1) at (-4.2,0.05) {Szenario \&\\Parameter};
\node[comp, fill=blue!14] (ui2) at (0,0.05)    {Laufsteuerung\\(Start/Pause/Schritt)};
\node[comp, fill=blue!14] (ui3) at (4.2,0.05)  {Vergleich\\\& Navigation};

% Row 2
\node[comp, fill=green!14] (app1) at (-4.2,-3.95) {Szenario-\\logik};
\node[comp, fill=green!14] (app2) at (0,-3.95)    {Simulationskern\\(RR, LCFS, Priority, MLFQ)};
\node[comp, fill=green!14] (app3) at (4.2,-3.95)  {Run-/Snapshot-\\Verwaltung};

% Row 3
\node[comp, fill=orange!18] (viz1) at (-4.2,-7.95) {Gantt-Zeitachse};
\node[comp, fill=orange!18] (viz2) at (0,-7.95)    {Warteschlange\\/ Statusansicht};
\node[comp, fill=orange!18] (viz3) at (4.2,-7.95)  {Kennzahlen\\\& Vergleich};

% Vertical main flow
\draw[flow] (ui1.south) -- (app1.north);
\draw[flow] (ui2.south) -- (app2.north);
\draw[flow] (ui3.south) -- (app3.north);

\draw[flow] (app1.south) -- (viz1.north);
\draw[flow] (app2.south) -- (viz2.north);
\draw[flow] (app3.south) -- (viz3.north);

% Internal logic flow WITHOUT labels (for readability)
\draw[flow] (app1.east) -- (app2.west);
\draw[flow] (app2.east) -- (app3.west);

% Feedback
\draw[feedback] (viz2.north) .. controls +(0,1.3) and +(0,-1.3) .. (ui2.south);

% Tiny legend
\node[anchor=west, font=\scriptsize] at (-6.45,-9.75)
{Hauptdatenfluss: oben $\rightarrow$ unten \quad gestrichelt: interaktive Rückkopplung};
\end{tikzpicture}
    \caption{Komponentenmodell der Anwendung mit den Ebenen Benutzeroberfläche, Anwendungslogik und Visualisierung.}
    \label{fig:komponentenmodell}
\end{figure}

Die Pfeile in Abbildung~\ref{fig:komponentenmodell} repräsentieren den Datenfluss vom Eingabeszenario über die Simulationsberechnung zur Visualisierung. Steueraktionen lösen keine fachlich getrennte Parallelwelt aus, sondern wirken auf denselben konsistenten Ablaufzustand.

% Placeholder: Scheduling Visualizer Übersichtsdiagramm
\begin{figure}[htbp]
    \centering
    \fbox{%
        \begin{minipage}[c][8cm][c]{14cm}%
            \centering
            \textit{[Scheduling Visualizer -- Übersichtsdiagramm: Screenshot der Hauptoberfläche mit Fokusansicht, Vergleichsübersicht und Kennzahlen]}
        \end{minipage}%
    }%
    \caption{Übersicht der Benutzeroberfläche: Eingabebereich, Simulationsfenster und Metriken-Anzeige.}
    \label{fig:app-overview}
\end{figure}


\section{Komponenten und Verantwortlichkeiten}

\subsection{Szenario- und Konfigurationskomponente}

Diese Komponente verwaltet die Eingabedaten der Prozesse (z.\,B. Ankunftszeit, Burst-Zeit, Priorität) sowie algorithmusspezifische Parameter (z.\,B. Zeitquantum). Zusätzlich verantwortet sie Validierung und Normalisierung der Eingaben, sodass alle Läufe auf strukturell konsistenten Daten aufsetzen.

\subsection{Simulationskomponente}

Die Simulationskomponente bildet den fachlichen Kern. Sie erzeugt aus Szenario und Verfahrensdefinition den vollständigen Ablauf als deterministische Folge von Zustandsübergängen. Ergebnisse werden nicht nur als Endzustand abgelegt, sondern in auswertbaren Strukturen bereitgestellt (Segmente, Ereignisse, Kennzahlen).

\subsection{Visualisierungs- und Interaktionskomponente}

Diese Komponente übersetzt Simulationsdaten in verständliche Ansichten: zeitlicher Ablauf (Gantt), Warteschlangen-/Statusdarstellung und Kennzahlen. Gleichzeitig steuert sie die Wiedergabe über Interaktionsfunktionen und ermöglicht den Vergleich mehrerer Läufe auf identischer Datenbasis.

\section{Datenfluss und Zustandsmodell}

Der Datenfluss folgt einem eindeutigen Muster:

\begin{enumerate}
    \item Eingabeszenario wird erfasst und validiert.
    \item Simulationskern berechnet den Ablauf für ein Verfahren.
    \item Ergebnisstrukturen (Ereignisse, Segmente, Snapshots, Metriken) werden bereitgestellt.
    \item Darstellungsschicht rendert die aktuelle Sicht und reagiert auf Interaktionsbefehle.
\end{enumerate}

Entscheidend ist, dass Visualisierung und Kennzahlen auf denselben Ergebnisstrukturen basieren. Dadurch entstehen keine widersprüchlichen Anzeigen zwischen Ablaufdarstellung und numerischer Bewertung.

\section{Interaktionsmodell für den Lernprozess}

Das Interaktionsmodell folgt einem didaktisch motivierten Nutzungsfluss:

\begin{enumerate}
    \item Szenario definieren,
    \item Verfahren auswählen,
    \item Ablauf steuern und beobachten,
    \item Ergebnisse interpretieren,
    \item Verfahren vergleichen.
\end{enumerate}

Die Interaktion ist damit kein rein technischer Bedienvorgang, sondern Teil des Lernprozesses. Die Steuerung beeinflusst die Sicht auf den Ablauf (z.\,B. Schrittmodus), ohne die fachliche Berechnung selbst zu verändern.

\section{HCI-Leitlinien im Entwurf}

Für eine Lernanwendung ist die UI ein fachlich relevantes Element. Der Entwurf folgt daher grundlegenden HCI-Leitlinien:

\begin{itemize}
    \item zentrale Zustände müssen direkt sichtbar sein,
    \item wichtige Ereignisse (Präemption, Kontextwechsel, Abschluss) müssen klar markiert werden,
    \item Bedienung und Analyse sollen in einem konsistenten Ablauf zusammenwirken,
    \item die Oberfläche soll Orientierung bieten und nicht zusätzliche kognitive Last erzeugen.
\end{itemize}

Damit wird sichergestellt, dass die Anwendung nicht nur technisch korrekt arbeitet, sondern auch als didaktisches Werkzeug nutzbar ist.

\section{Erweiterbarkeit und Entwurfsgrenzen}

Die Architektur ist auf Erweiterbarkeit ausgelegt: Neue Scheduling-Verfahren können als zusätzliche Strategien im Simulationskern ergänzt werden, ohne die UI grundlegend umzubauen. Ebenso lassen sich neue Kennzahlen oder Ansichten integrieren, sofern sie auf den vorhandenen Ergebnisstrukturen aufsetzen.

Eine bewusste Entwurfsgrenze liegt in der Fokussierung auf präemptive Verfahren und den Lehrkontext. Die Architektur ist zwar übertragbar, aber nicht als vollständiger Ersatz für produktive Betriebssystem-Scheduler konzipiert.